\section{Discussion}\label{sec:discussion}
\begin{CJK*}{UTF8}{gbsn}

\nextrevision{本文结果表明，图谱知识能否改善零样本音频分类，不仅取决于是否提供了类别名称之外的信息，还与关系选择、知识表达及分数融合方式有关。SAKI分别从这三个环节组织知识增强过程，完整配置取得最高的五数据集等权平均Hit@1和MRR，支持三个模块在总体表现上的互补作用。这些结果说明，在补充图谱知识的同时，还需要考虑如何针对当前音频选择关系、表达关系语义，以及协调知识证据与原始CLAP预测。}

\nextrevision{AAKV的作用需要分别从语言化对照和组合消融理解。在相同的关系选择与融合配置下，AAKV的Hit@1在五个测试集上均高于直接拼接和原始三元组表示，支持其在该配置下作为知识文本表示方式的有效性。AAKV保留三元组中的实体和关系方向，并将其组织为面向声音的文本，为CLAP利用关系语义提供了更明确的表达。另一方面，组合消融显示，AAKV单独启用时平均增益有限，且部分数据集的性能下降；与其他模块共同使用后，完整配置取得最高的五数据集等权平均Hit@1和MRR。这些结果表明，AAKV在不同模块配置下的收益存在差异，其单独使用的效果不足以概括其在完整方法中的贡献。}

\nextrevision{关系选择对照进一步支持了样本级关系选择的设计。在相同知识表示和融合配置下，该策略优于所比较的固定、高频和随机关系策略；所选关系组合的多样性也表明，选择结果随输入音频发生变化。前者支持其分类有效性，后者体现其样本依赖性。iKnow-audio采用经过筛选的关系集合检索尾实体，以扩展类别提示\,\cite{olvera-etal-2025-iknow}；SAKI则进一步依据当前音频下关系专家的预测共识与分类间隔，执行样本级关系选择，使参与知识增强的关系能够随输入调整。}

\nextrevision{不同评价指标的变化揭示了知识增强的收益边界。FSD50K上，在AAKV与样本级关系选择的组合中加入分层融合后，Hit@1基本持平，MRR略有提高；TUT2017上，完整方法相较于iKnow\textsuperscript{\dag}提高了Hit@1和MRR，但Hit@5略有下降。本文主要以Hit@1和MRR评价首位分类表现与真实类别的排序质量，并通过Hit@3和Hit@5补充考察前若干候选类别对真实标签的覆盖情况。因此，上述结果支持完整方法在本文主要评价指标上的总体优势，同时表明这种优势并未同步体现在所有评价指标上。逐样本分析同样观察到错误纠正、排名改善和预测退化，且五个测试集上的正向预测转移数量均超过负向转移数量，说明总体性能提升体现为逐样本排序变化带来的净改善。}

\nextrevision{SAKI仍受到知识质量、基础模型和计算成本的约束。AAKV能够约束三元组的文本表达，但不能校正三元组本身的错误；关系选择与知识评分也依赖CLAP对音频和文本的匹配能力。此外，尽管尾实体检索、AAKV文本生成和文本编码均在离线阶段完成并缓存，在线关系专家评估仍带来额外开销。当前测量中，SAKI的平均在线延迟为86.52 ms/样本，是CLAP的3.55倍，说明性能提升伴随着计算和存储成本的增加。当前结论基于固定的CLAP模型、RotatE评分函数及候选关系池；后续可考察证据可信度评估能否减少预测退化、轻量级关系预筛选能否降低在线开销，并在其他音频--语言模型与知识图谱上验证方法的适用范围。}

\end{CJK*}
